\section*{Chapter \ref{ch:iv1:games} --- Games}
\markboth{\protect\omTocMark{}{Answers: Games}}{\protect\omTocMark{}{Answers: Games}}

\begin{solution}{iq:iv1:games:1}
Go first and take 6, leaving 91. If taking the last stone loses, go first and take 5, leaving 92.

Multiples of 7 are losing for the player to move\index{losing positions}. From a multiple of 7, any move leaves a non-multiple; from a non-multiple, you can always get to a multiple. So whatever the opponent takes, you take the rest of 7. Since $97 = 13 \cdot 7 + 6$, you take 6.

If the last stone loses, you want the opponent to face exactly 1 stone. So the losing positions are one above a multiple of 7, and you need to leave 92.

\iqwrong{Saying that the losing total is the same in both versions. The target shifts by one, because the last stone is now a stone you want to avoid.}
\end{solution}

\begin{solution}{iq:iv1:games:2}
The first player wins by taking 3 stones from the larger pile, which makes both piles 5. If the piles start equal, the second player wins.

Equal piles lose for the player to move\index{mirror strategy}. Whatever that player takes from one pile, the opponent takes the same from the other, so the piles stay equal and the opponent gets the last stone. Unequal piles win, because one move makes them equal. A brute-force table up to piles of 9 confirms the rule.

\iqwrong{Thinking that the larger pile is an advantage in itself, and trying to keep it. Only the difference matters, and the winner is the one who hands over a difference of 0.}
\end{solution}

\begin{solution}{iq:iv1:games:3}
Each player shows heads with chance $3/8$, and the game is worth $-\$1/8$ to you per round.

A mix must make the other player indifferent\index{mixed strategy}, or that player would exploit it. If I show heads with chance $q$, your two choices earn
\[
3q - 2(1 - q) \quad\text{and}\quad -2q + (1 - q),
\]
and these are equal at $q = 3/8$, which gives a value of $-1/8$. Your mix $p$ solves the same equations, and it is also $3/8$.

As a second route, add up the four outcomes under those mixes: $\tfrac{9}{64} \cdot 3 + \tfrac{25}{64} \cdot 1 - \tfrac{30}{64} \cdot 2 = -\tfrac18$.

The four payoffs add up to zero, and so the game looks fair, but it is not. The big payoff is the one you can rarely collect: heads against heads happens only $9/64$ of the time, because I avoid it.

\iqwrong{Saying each should show heads more often, since the big win comes from heads. If you show heads more, I would show tails, and you would pay \$2 more often. Your payoffs set my mix, not yours.}
\end{solution}

\begin{solution}{iq:iv1:games:4}
Fully rational players write 0, or 1, since with whole numbers both are stable. In a real room you write a number above that, because other people do not go all the way down.

Two-thirds of the average can never exceed $66.7$, so any number above 67 does no better than 67, whatever the others write\index{iterated dominance}. If everybody knows that, the average is at most 67, so nothing above 45 is worth writing. Repeating the argument brings the top down to 1. If all write 0, the target is 0. If all write 1, the target is $\tfrac23$. In a room of five or more, it stays closer to 1 than to 0 even after one person moves to 0, so nobody gains by moving down.

In a real room, the question is how many steps others take. If they pick at random, averaging 50, the best reply is 33. Thinking that they are all at 33, the reply is 22, and then 15. So pick a low number that is one step beyond what the others do, not the theoretical answer.

\iqwrong{Writing 0 and claiming victory. It is the equilibrium if everyone is rational and knows it, but the question asks you to win. Against real people, too much reasoning loses just like too little.}
\iqwhy{This is a model of markets where value depends on what others think that others think. The rational answer and the profitable answer can differ, and being one step ahead beats being right.}
\end{solution}

\begin{solution}{iq:iv1:games:5}
Play rock $1/2$, paper $1/6$ and scissors $1/3$ of the time. Rock is played most, even though a win with rock pays least.

In a symmetric zero-sum game like this one, the value is 0 and each choice must earn 0 against the mix. Let the mix be $(r, p, s)$\index{mixed strategy}:
\[
\text{rock: } s - 2p = 0, \quad \text{paper: } 2r - 3s = 0, \quad \text{scissors: } 3p - r = 0.
\]
So $r = 3p$ and $s = 2p$, and these add up to $6p = 1$.

It looks odd, because you would expect the highest-paying choice to be most common. But the frequencies are not set by what a choice earns for you. They are set by what they do to the other player. Rock is frequent because it has to deter the choice that pays most: if rock were rare, scissors would earn its \$3 with no risk. A linear program on the $3 \times 3$ matrix gives the same mix.

\iqwrong{Mixing in proportion to the payoffs, $1:2:3$. Against that mix, rock earns $\tfrac36 \cdot 1 - \tfrac26 \cdot 2 = -\tfrac16$, paper earns $\tfrac16 \cdot 2 - \tfrac36 \cdot 3 = -\tfrac76$, and scissors earns $\tfrac26 \cdot 3 - \tfrac16 \cdot 1 = \tfrac56$, so the opponent plays scissors every time.}
\end{solution}

\begin{solution}{iq:iv1:games:6}
Go first. The nim-sum is $3 \oplus 5 \oplus 7 = 1$, which is not zero. There are three winning moves: to $(2, 5, 7)$, $(3, 4, 7)$ and $(3, 5, 6)$.

The rule\index{Nim}\index{nim-sum}: a position loses for the player to move exactly when the bitwise exclusive-or of the piles is 0. Here is the proof, in three steps.

First, the empty position has nim-sum 0, and it is lost.

Second, any move from nim-sum 0 gives a non-zero nim-sum, because changing one pile changes the bits, and the rest stay fixed.

Third, from a non-zero nim-sum $s$, there is a move to 0. Take a pile with the top bit of $s$ set, and replace it by pile $\oplus\ s$. That new value is smaller, since the top bit of $s$ goes from 1 to 0, and it makes the total 0.

Here $s = 1$, whose top bit is bit 0. All three piles are odd, so any of them can be used: $3 \to 2$, $5 \to 4$, $7 \to 6$. A brute-force search of all piles up to 5, and of this position, confirms the rule and the three moves.

\iqwrong{Trying to reason by parity of the number of stones, or to keep the piles balanced by size. The binary digits decide, and piles of $1, 2, 3$ lose for the player to move even though they look uneven.}
\end{solution}

\begin{solution}{iq:iv1:games:7}
The Grundy values for 0 to 8 stones are $0, 1, 0, 1, 2, 3, 2, 0, 1$, with period 7. The first player wins with piles of 10 and 6, because $g(10) \oplus g(6) = 1 \oplus 2 = 3$. The winning moves are to take 4 from the 10, leaving 6 and 6, or to take 3 from the 6, leaving 10 and 3.

The Grundy value of a position is the smallest non-negative integer not among the values of its moves\index{Grundy value}. A position with value $v$ behaves in a sum like a Nim pile of size $v$\index{Sprague--Grundy theorem}. So for $n = 4$, the moves lead to 3, 1 and 0, with values $1, 1, 0$, so $g(4) = 2$. For $n = 5$ the moves lead to 4, 2, 1, with values $2, 0, 1$, so $g(5) = 3$.

A sum loses when the values' nim-sum is 0. You need a move in one pile that matches the value of the other. From 10, a value of 2 is reachable by taking 4. From 6, a value of 1 is reachable by taking 3, as $g(3) = 1 = g(10)$. A brute-force table over pairs of piles finds these two moves and no others.

\iqwrong{Using plain Nim on the pile sizes, $10 \oplus 6 = 12$. Pile sizes only work if any number of stones can be removed. Here you must use the Grundy values.}
\end{solution}

\begin{solution}{iq:iv1:games:8}
A proposes 98 for A, 0 for B, 1 for C, 0 for D and 1 for E. If a proposal needs more than half, A proposes 97 for A, 1 for C and 2 for one of D or E.

Work backwards\index{backward induction}\index{pirate game}, with at least half needed.

With D and E, D votes for D's own plan, which is half: $(100, 0)$.

With C, D and E: C needs one vote. E gets 0 otherwise, so E takes 1: $(99, 0, 1)$.

With B, C, D and E: B needs one more vote. D gets 0 otherwise: $(99, 0, 1, 0)$.

With all five: A needs two more votes, from those who would get 0 under B, and they are C and E: $(98, 0, 1, 0, 1)$.

With more than half needed, D is doomed when only D and E are left, since E can reject and take all. So D will back C for free, and C keeps 100 with D's vote. B needs three of the four votes, and D and E each get 1, for $(98, 0, 1, 1)$. A needs three of five, and the cheapest are C, for 1, and one of D or E, for 2.

\iqwrong{Giving each pirate an equal share to get votes. A voter compares with the next round, not with the others, and a vote costs one coin above what that voter would get next round.}
\iqwhy{In any vote or negotiation the price of a voter depends on what that voter gets if the current offer fails. That is the logic of a bargaining position.}
\end{solution}

\begin{solution}{iq:iv1:games:9}
With three rounds, A gets $0.91$. With no limit, A gets $1/(1 + 0.9) = 10/19 \approx 0.526$.

For three rounds, work backwards\index{backward induction}, in the money of that round. In round 3, A offers B nothing and keeps 1, because B gets nothing by rejecting. In round 2, B must offer A what round 3 is worth to A today, which is $0.9$, and B keeps $0.1$. In round 1, A must offer B $0.9 \cdot 0.1 = 0.09$, and A keeps $0.91$.

With no limit, let $x$ be what the proposer gets\index{alternating-offer bargaining}. The game looks the same after a rejection, so the responder must be given $\delta x$, and the proposer keeps $1 - \delta x$. Then $x = 1 - \delta x$, which gives $x = \frac{1}{1 + \delta}$. Iterating the backward induction over many rounds converges to this.

The limits make sense. At $\delta = 0$, the proposer takes all, and as $\delta \to 1$ both get a half. Proposing first is worth more when delay is costly.

\iqwrong{Using the last-mover result for the unlimited game. The large share of $0.91$ comes from A having the final offer, and with no last round, that advantage is gone.}
\end{solution}

\begin{solution}{iq:iv1:games:10}
The taker aims left with chance $8/15$, the keeper dives left with chance $3/5$, and the chance of a goal is $0.74$. After the improvement the taker aims right less often: $7/16$ instead of $7/15$.

The taker's mix leaves the keeper indifferent\index{mixed strategy}. With $p$ the chance of aiming left, the goal chance is
\[
0.6p + 0.9(1 - p) \text{ if the keeper dives left, and } 0.95p + 0.5(1 - p) \text{ if right.}
\]
These agree at $p = 8/15$, at a value of $0.74$. The keeper's chance $q$ of diving left solves $0.6q + 0.95(1 - q) = 0.9q + 0.5(1 - q)$, so $q = 3/5$.

With the $0.9$ raised to $0.95$, the first expression changes, and $p = 9/16$. So aiming right falls from $7/15 \approx 0.467$ to $7/16 \approx 0.438$, and the value rises to about $0.753$.

The surprise is correct. The taker's mix has to keep the keeper indifferent, and when aiming right has become more dangerous, the keeper dives left less. That keeper's shift makes it better to aim left. The improved option is used less, and the taker still gains. A linear program gives both results.

\iqwrong{Saying that you should use your improved skill more. In a game with a responding opponent, a change in your payoffs moves the opponent's mix, and yours is set by what makes the opponent indifferent.}
\iqwhy{An improved signal or trading venue does not always mean you should use it more. If other players respond, the result is in the equilibrium, not in your own table of payoffs.}
\end{solution}

\begin{solution}{iq:iv1:games:11}
In a second-price auction, bid your value. In a first-price auction, bid $\frac{n - 1}{n}v$. In both, the seller expects $\frac{n - 1}{n + 1}$.

Second price: your bid decides only whether you win, never what you pay\index{second-price auction}. If the highest other bid is below your value, you want to win, and if it is above, you do not. Bidding your value does both, so it is dominant.

First price: suppose the others bid $\beta v$\index{first-price auction}. Bidding $b$ wins when all $n - 1$ others have values below $b/\beta$, so the payoff is $(v - b)(b/\beta)^{n - 1}$. Setting the derivative to zero gives $b = \frac{n - 1}{n}v$, for any $\beta$, and so this is consistent. A scan over bids at $v = 0.9$ confirms the maximum for $n = 2, 3, 5$.

Revenue: second-price pays the second-highest value, whose mean is $\frac{n - 1}{n + 1}$. First-price pays $\frac{n - 1}{n}$ times the highest value, whose mean is $\frac{n}{n + 1}$, which gives the same. This is revenue equivalence\index{revenue equivalence}, and it is an equality of means only. The first-price revenue has a lower variance in the simulation.

\iqwrong{Assuming a first-price auction earns more because the winner pays their own bid. Bidders shade their bids to allow for that, and in this setup the shade exactly cancels the difference.}
\end{solution}

\begin{solution}{iq:iv1:games:12}
Each does it with probability $3/4$ for $n = 2$ and $1/2$ for $n = 3$. Nobody does it with chance $1/16$ and $1/8$, and this rises towards $1/4$ as $n$ grows, while each person's own probability falls to 0.

A mixed strategy\index{mixed strategy} is used only if both choices pay the same. Doing it pays 3 for sure. Not doing it pays 4 if at least one of the other $n - 1$ does it. With probability $p$ for each, the condition\index{volunteer's dilemma} is
\[
3 = 4\bigl(1 - (1 - p)^{n - 1}\bigr), \qquad\text{that is}\qquad (1 - p)^{n - 1} = \tfrac14.
\]
So $p = 3/4$ for $n = 2$ and $p = 1/2$ for $n = 3$. The chance that nobody does it is $(1 - p)^n = (\tfrac14)^{n/(n - 1)}$, which gives $1/16$, $1/8$, and tends to $1/4$.

Each person's payoff is 3 whatever $n$ is, since doing it is always one of their best choices. So a larger group helps nobody here.

\iqwrong{Saying that with more people, the task is more likely to get done. Each person relies more on the others as the group grows, and this more than offsets the extra hands. Exactly one person doing it is also an equilibrium, but only if the group can agree who.}
\iqwhy{A risk that anyone on a desk could flag, at a small private cost, has this shape. A larger group is not a safer one unless someone is named as responsible.}
\end{solution}

\begin{solution}{iq:iv1:games:13}
The first player wins, by a strategy-stealing argument\index{strategy stealing}. On a square bar, the explicit winning move eats the square diagonally next to the bottom-left one, which leaves two equal arms; after that, the first player mirrors.

Suppose that the first player eats only the top-right square. If that leaves a losing position for the opponent, the first player has won. Otherwise, the second player has a winning reply, which eats some square $(i, j)$. But eating $(i, j)$ also eats the top-right square, as that lies above and to the right of every square. So the first player could have played $(i, j)$ straight away and reached the same position. Either way, a winning first move exists. It is non-constructive, and it tells you nothing about the move.

Square: eat the square at the diagonal neighbour of the corner. What is left is a bottom row and a left column of the same length. Whatever the opponent eats from one arm, eat the same from the other. The opponent is left to eat the corner.

A brute-force game tree for all bars up to 5 by 5 confirms that the first player wins every bar except the $1 \times 1$.

\iqwrong{Trying to find a pattern for the winning move in general. No simple formula for it is needed here, and the question is about existence. You should say so, and give the argument.}
\end{solution}

\begin{solution}{iq:iv1:games:14}
With 201 the proposer survives with no coins, and with 202 the same. With 203 the proposer dies, and with 204 the proposer survives. Beyond that, the proposer survives only at $200 + 2^k$ pirates: 208, 216, 232, 264, 328, and so on. With 300 pirates, 36 are thrown overboard.

Coins can buy at most 100 votes\index{pirate game}. With 200 pirates, the proposer needs 100 votes, pays 99 and keeps 1. With 201, 101 votes are needed: 100 bought, plus the proposer's own, leaving none. With 202, 101 votes are still needed, and the same plan works.

With 203, 102 votes are needed, and there are only 101 votes available. So the proposer dies, whatever the plan. This creates a new source of votes: a proposer who would die in turn will vote for anything, in order to live.

With 204, 102 votes are needed: 100 bought, the proposer, and the next one, who would die otherwise.

In general, let $m$ be the last size at which the proposer survives. At size $n = m + j$, the proposer has $j$ self-preserving votes, including their own, plus 100 bought. They need $\lceil n/2 \rceil$, which gives $j \ge m - 200$. So from $m = 202$ you get 204, then 208, 216 and so on, and each survivor has twice the distance above 200.

For 300, the next survivor below is 264, so the 36 top pirates die one after another. At 299 votes needed are 150, and there are $35 + 100$ votes available, and it stays short all the way down to 264.

\iqwrong{Extending the coin allocation in a simple way past 200, or thinking that everyone dies once coins run out. Fear of being next creates free votes, which is why some of the larger sizes work.}
\end{solution}

\begin{solution}{iq:iv1:games:15}
X bets the ace always, checks the king, and bluffs the queen with chance $1/3$. Y calls with the ace, calls with the king with chance $1/3$, and folds the queen. The game is worth $1/18$ of a dollar to X, about $0.056$.

Some choices are dominated\index{bluffing frequency}. With the ace, X should bet, and Y should call. With the queen, Y must fold. With the king, betting is bad for X: the queen folds, which is what checking would win anyway, and the ace calls, which costs more. So only two numbers are left.

Bluff frequency $b$ makes a king indifferent. When Y holds the king and faces a bet, X has the ace (1 way) or a bluffing queen ($b$ ways). Calling earns $\frac{-2 + 2b}{1 + b}$, and folding earns $-1$. These agree at $b = 1/3$.

Calling frequency $c$ makes a queen indifferent. When X holds the queen, checking loses 1. Betting against the king earns $1 - c$ when Y folds and $-2c$ when Y calls, and betting against the ace loses 2. So $-1 = \tfrac12\bigl((1 - c) - 2c\bigr) + \tfrac12(-2)$, which gives $c = 1/3$.

The value comes from a linear program on the $8 \times 8$ table of pure strategies, $1/18$, and a simulation of the two mixes matches this. X earns it because X alone can decide whether more money goes in, so Y has to pay for the right to catch a bluff. The small edge comes from this choice, and not from any better cards, since the deal is fair to both players.

\iqwrong{Bluffing half of the time, or never. If you never bluff, the king always folds, and your ace earns less than it should. If you bluff too much, the king always calls. The correct rate makes the caller indifferent. A call costs one dollar to win a pot of three, so one bluff for every three value bets.}
\end{solution}

\begin{solution}{iq:iv1:games:16}
The best reserve is $1/2$, with revenue $5/12$, against $1/3$ with no reserve. It does not depend on the number of bidders. A reserve lowers the total surplus, from $2/3$ to $7/12$.

Bidding your value is still dominant. If both values are above $r$, the price is the lower one. If only one is, the price is $r$. If none is, no sale happens\index{reserve price}. So
\[
R(r) = \int_r^1 x \cdot 2(1 - x)\,dx + r \cdot 2r(1 - r) = \frac13 + r^2 - \frac43 r^3.
\]
Setting $R'(r) = 2r - 4r^2 = 0$ gives $r = \tfrac12$, and $R(\tfrac12) = 5/12$. A simulation matches this.

Why does it not depend on $n$? The reserve is a rule for a single sale. It matters only when one bidder is above it and all others below, and then you are in effect setting a price for one buyer with a uniform value. A buyer whose value is above $r$ pays $r$ with chance $1 - r$, whose product is maximised at $\tfrac12$, whoever else is there. With 4 bidders, a scan of reserves gives the same result.

For surplus, the good goes to the highest bidder unless both values are below $\tfrac12$, which happens with chance $\tfrac14$. The surplus is then $\int_{1/2}^1 x \cdot 2x\,dx = 7/12$, instead of $E[\max] = 2/3$.

\iqwrong{Saying that a reserve can only hurt because it risks no sale. It costs $\tfrac14$ of the sales, but it raises the price paid in the other cases, and the seller comes out ahead.}
\end{solution}

\begin{solution}{iq:iv1:games:17}
The first player wins, with one move only: take 2 coins from step 1, leaving $(1, 1, 4, 1, 5)$. Only the odd steps matter, and their nim-sum is $3 \oplus 4 \oplus 5 = 2$.

Treat the odd steps as Nim piles\index{staircase Nim}. Moving coins from an odd step to the step below is a Nim move on that pile, as the coins land on an even step, which does not count. Moving coins from an even step lands them on an odd step, which enlarges a pile. But the opponent can always cancel that by moving the same coins on, from that odd step to the even step below (from step 1 they leave the staircase, and that counts as the same thing). So a move from an even step gains nothing. The game must end, because the coins go only downwards.

So the player facing a nim-sum of 0 on the odd steps loses. Here you need to make it 0. With $s = 2$, the pile that has bit 1 set is 3, which becomes $3 \oplus 2 = 1$. So you take 2 from step 1. Moves from even steps do not help: they make the sums 5 or 3.

A brute-force search over all moves from this position finds that single move, and the rule matches the brute force over all small staircases.

\iqwrong{Using all five piles in the nim-sum. The even steps can always be undone, so they are not real choices.}
\end{solution}

\begin{solution}{iq:iv1:games:18}
The losing positions are $(0, 0)$, $(1, 2)$, $(3, 5)$, $(4, 7)$, $(6, 10)$, $(8, 13)$ and $(9, 15)$. The position $(9, 15)$ loses for the player to move, and $(10, 15)$ wins, by taking 1 from the first pile. The closed form is $\bigl(\lfloor k\varphi \rfloor, \lfloor k\varphi^2 \rfloor\bigr)$, with $\varphi = (1 + \sqrt{5})/2$\index{Wythoff's game}.

Build the list by hand. A losing position cannot share a pile size with another losing position, since one pile move would connect them. And it cannot share a difference with another, since a move on both piles keeps the difference. So at step $k$, take the smallest positive integer not yet used as the smaller pile, and add $k$ to get the larger. Every other position can reach a losing position by one move, as a brute-force table up to 40 confirms.

The sequences $\lfloor k\varphi \rfloor$ and $\lfloor k\varphi^2 \rfloor$ partition the positive integers exactly, by Beatty's theorem\index{Beatty's theorem}, since $1/\varphi + 1/\varphi^2 = 1$. Also, $\varphi^2 = \varphi + 1$, so $\lfloor k\varphi^2 \rfloor = \lfloor k\varphi \rfloor + k$, which is the required difference. So the formula fits the building rule.

\iqwrong{Using nim-sum again. With diagonal moves allowed, this is not a sum of independent piles, and $9 \oplus 15 = 6$ is not zero, even though the position is losing.}
\end{solution}

\begin{solution}{iq:iv1:games:19}
The game is worth $-1/18$ to Player 1, so the second player has the edge. There is a whole family of equilibria for Player 1. For $\alpha \in [0, \tfrac13]$, bet the jack with chance $\alpha$, always check the queen and call a bet with chance $\tfrac13 + \alpha$, and bet the king with chance $3\alpha$, otherwise checking and calling. Player 2's equilibrium strategy is: with the king, always bet or call; with the queen, check when checked to, and call with chance $\tfrac13$; with the jack, bet with chance $\tfrac13$ after a check, and fold to a bet.

Many choices are dominated\index{Kuhn poker}. The king never folds, the jack never calls, and the queen is too strong to bluff with and too weak to bet for value. So the jack bluffs, and the king sometimes traps.

Now use indifference\index{bluffing frequency}. Player 2 bluffs the jack with chance $b$ after a check. Player 1 with the queen faces that bet from a king (always) or a jack ($b$). Calling earns $\frac{-2 + 2b}{1 + b}$, and folding earns $-1$, so $b = \tfrac13$.

Player 2's queen calling a bet with chance $c$ must keep Player 1's jack indifferent. Checking with the jack loses 1 against both. Betting earns $\tfrac12\bigl((1 - c) - 2c\bigr) + \tfrac12(-2)$. So $c = \tfrac13$.

Player 2's queen is indifferent to a bet when the bets are three-quarters kings. That is why Player 1's king bets with chance $3\alpha$ when the jack bluffs with chance $\alpha$. Player 1's queen call rate of $\tfrac13 + \alpha$ keeps Player 2's jack indifferent.

A linear program on the $27 \times 64$ table of pure strategies gives $-1/18$. Each member of the family, tested at $\alpha = 0$, $0.1$, $0.2$ and $\tfrac13$, earns at least that against all 64 replies, and Player 2's strategy holds all 27 pure strategies to at most that.

\iqwrong{Assuming the first to bet has the advantage. Acting second is worth more here, because you see what the other player does before you decide. Also, assuming that a unique equilibrium exists, when Player 1 has a whole line of them with the same value.}
\iqwhy{An edge of $1/18$ of an ante comes entirely from position, with both players making no mistakes. Most of what you can earn in practice comes from an opponent's departure from these frequencies.}
\end{solution}

\begin{solution}{iq:iv1:games:20}
X bets with hands below $0.1$ and above $0.7$, and checks in between. Y calls with hands above $0.4$. The game is worth $\$0.10$ to X.

Let Y call above $c$\index{continuous poker}\index{bluffing frequency}. Checking with hand $x$ earns $2x - 1$. Betting earns $c$ from Y's folds. With $x$ below $c$, every call loses 2, so betting earns $c - 2(1 - c) = 3c - 2$, whatever $x$ is. With $x$ above $c$, betting earns $c + 2\bigl((x - c) - (1 - x)\bigr)$.

Three conditions determine the solution:
\begin{itemize}
\item A bluffing hand at $a$ is indifferent: $2a - 1 = 3c - 2$.
\item A value hand at $b$ is indifferent: $2b - 1 = c + 2(2b - c - 1)$, so that $b = (1 + c)/2$.
\item Y at $c$ is indifferent between folding, which loses 1, and calling, which wins 2 against a bluff and loses 2 against a value hand: $\frac{2a - 2(1 - b)}{a + 1 - b} = -1$, so that $3a = 1 - b$.
\end{itemize}
Solving gives $c = 0.4$, $a = 0.1$ and $b = 0.7$. Betting with the very worst hands works, because they lose anyway when they check, and they are the hands whose showdown value is lowest.

For a bet size $B$, the same steps give $a = \frac{B}{(B + 1)(B + 4)}$, $c = \frac{B(B + 3)}{(B + 1)(B + 4)}$ and $b = \frac{B^2 + 4B + 2}{(B + 1)(B + 4)}$. The value of the game equals $a$, and it is largest at $B = 2$, where it is $1/9$. Numerical integration confirms the value at $B = 0.5$, 1, 2 and 4, and a simulation of the rules at $B = 1$ gives $0.10$. For the given range, every other calling threshold gives X at least as much, which confirms that $c$ is Y's best reply.

\iqwrong{Bluffing with middling hands, because those are nearest to winning. Bluffing has a constant payoff, so you bluff with the hands that have the least to lose by not checking, which are the very worst. Middling hands prefer to check and win at a showdown when they can.}
\iqwhy{The ratio of bluffs to value bets here is $0.1:0.3$. A call costs \$1 to win a pot of \$3, and one bluff for every three value bets makes the caller indifferent. Any quoted or displayed size has to be readable in the same way, as a mix of strength and bluff.}
\end{solution}
